\chapter{Tables and Figures}\label{chap:tables-figures}

This chapter covers the first half of the fourth block:
tables, how to draw their lines, floats and captions, and images.
Tables come first, because figures work the same way and are then nearly free.

\section{Tables}\label{sec:tables}

A table is a \texttt{tabular} environment.
Its argument has one letter per column, \code|&| moves to the next column, and \code|\\| ends a row:

\slidecode{04-warmup.tex}
\sliderender{04-warmup.pdf}

Both characters are old friends:
\code|&| is the column separator from the reserved characters (\autoref{sec:reserved}),
and \code|\\| is one of the places where a forced line break belongs (\autoref{sec:spaces}).
With a header row and four columns:

\slidecode{04-tabular.tex}
\sliderender{04-tabular.pdf}

The table has no caption and no number yet, and it is set right where it stands, like a big letter in the line.
\autoref{sec:floats} puts it into a float.

\subsection{Column types}\label{sec:column-types}

\autoref{tab:column-types} lists the column types.
The last two come from the package \texttt{array}, which many other packages load anyway.

\begin{table}[htbp]
    \centering
    \caption{Column types of \texttt{tabular}.}
    \label{tab:column-types}
    \begin{tabular}{@{}ll@{}}
        \toprule
        \textbf{Type} & \textbf{Column} \\
        \midrule
        \code|l|, \code|c|, \code|r| & left-aligned, centred, right-aligned \\
        \code|p{3cm}|                & fixed width, text wraps, top-aligned \\
        \code+|+                     & a vertical line between two columns \\
        \code|@{...}|                & replaces the space between two columns, e.g.\ \code|@{}| for none \\
        \code|m{3cm}|, \code|b{3cm}| & like \code|p|, but vertically centred or bottom-aligned (\texttt{array}) \\
        \code|>{...}|                & code before every cell of the column (\texttt{array}) \\
        \bottomrule
    \end{tabular}
\end{table}

Text in an \code|l|, \code|c| or \code|r| column never wraps,
so a long text runs off the page.
For such a column, give it a width with \code|p|:

\owncode{ch06-pcolumn.tex}
\ownrender{ch06-pcolumn.pdf}

A \code|p| column is justified like a paragraph, which in a narrow column leaves wide gaps between the words.
\code|>{\raggedright\arraybackslash}p{5cm}| sets it ragged right instead
(\code|\arraybackslash| keeps \code|\\| working at the end of the row).
To make a table exactly as wide as the text, use \texttt{tabularx} (\autoref{sec:tabularx}).

Numbers in a column are easiest to compare right-aligned,
or aligned at the decimal point with the \code|S| column of \texttt{siunitx} (\autoref{sec:siunitx}).

\subsection{Cells over several columns}

\code|\multicolumn{2}{c}{...}| joins two cells of a row into one, centred.
Together with a partial rule (\autoref{sec:table-lines}), it groups the columns under a common heading:

\owncode{ch06-multicolumn.tex}
\ownrender{ch06-multicolumn.pdf}

\section{Table lines}\label{sec:table-lines}

Lines are drawn with \code+|+ in the column types for vertical ones, and \code|\hline| for horizontal ones:

\slidecode{04-grid.tex}
\sliderender{04-grid.pdf}

The package \texttt{booktabs} offers a better way.
It adds three rules of different weight, \code|\toprule|, \code|\midrule| and \code|\bottomrule|:

\slidecodetitled{04-booktabs-preamble.tex}{Preamble}
\slidecodetitled{04-booktabs.tex}{Document}

Side by side, the grid on the left, \texttt{booktabs} on the right:

\sliderender{04-lines.pdf}

The \texttt{booktabs} table is how tables are set in books and journals.
Its manual \parencite{fear2020} gives two rules:
never use vertical lines, and never use double lines.
White space separates the columns well enough,
and the extra lines only make the table harder to scan.
Horizontal lines are used sparingly too: one above, one below the header, one at the end.
For more structure, there are
\code|\cmidrule(lr){2-3}|, a thinner rule under columns 2 to 3 (shortened at the left and right end),
and \code|\addlinespace|, a little extra space between groups of rows.

Without \code|\usepackage{booktabs}|, the three rules give the error \texttt{Undefined control sequence}:
they are not part of \hologo{LaTeX} itself.

\section{Floats}\label{sec:floats}

A large table or figure does not fit just anywhere:
in the middle of a paragraph it would leave a gap at the bottom of the page,
or be split across two pages.
So it goes into a \emph{float}, an environment that \hologo{LaTeX} may move to a better place:

\slidecode{04-table-float.tex}

\texttt{table} is the float, \texttt{tabular} the grid of cells inside it;
the two names are easy to mix up.
The float for images is \texttt{figure} (\autoref{sec:images}).

\subsection{Placement}

The optional argument lists where the float may go:

\begin{description}
    \item[\texttt{h}] here, at the place in the text where it is written
    \item[\texttt{t}] at the top of a page
    \item[\texttt{b}] at the bottom of a page
    \item[\texttt{p}] on a page of its own, with only floats on it
\end{description}

\code|[htbp]| allows all four.
Without an argument, the standard classes use \code|[tbp]|.

\hologo{LaTeX} tries the places of the list for the current page,
and if none of them fits, the float waits for the next page.
Roughly, a place fits when the float, together with the other floats already there, stays within some limits:
floats may take at most \qty{70}{\percent} of a page at the top and \qty{30}{\percent} at the bottom,
at least \qty{20}{\percent} of a page must remain text,
and a page of floats only (\texttt{p}) must be at least half full.
Floats of the same kind also keep their order:
a table that has to wait holds back every later table.
At \code|\clearpage|, and at the end of the document, all waiting floats are set.

So a table that moves to the next page is not a bug.
It is what floats are for,
and the reason they get a number: the text says \enquote{Table~1}, not \enquote{the table below}.
Fighting float placement takes a lot of time and rarely helps.
Write the text, and look at the placement once at the end.
If something is still badly placed:

\begin{itemize}
    \item Check that the argument allows enough places.
        \code|[h]| alone is a common reason for floats that pile up at the end of a chapter:
        if \texttt{h} does not fit, the float has hardly anywhere else to go.
        \hologo{LaTeX} notices this and changes \code|[h]| to \code|[ht]| by itself,
        with the warning \texttt{`h' float specifier changed to `ht'},
        but a float that fits neither place still waits for the next page.
    \item \code|!| in front, e.g.\ \code|[!htbp]|, relaxes the limits above for this float.
        It does not force anything.
    \item \code|\FloatBarrier| from the package \texttt{placeins} sets all waiting floats before it continues,
        without starting a new page.
        \code|\usepackage[section]{placeins}| puts one at every \code|\section|,
        so no float drifts into the next section.
    \item \code|[H]| from the package \texttt{float} forces the float to stay exactly where it is,
        so it is no longer a float at all.
        Use it only as a last resort:
        a float that does not fit leaves the rest of the page empty,
        and a document full of \code|[H]| is full of half-empty pages.
\end{itemize}

\subsection{Inside a float}

A float holds its content, a caption and a label, usually in this order (for figures, the caption comes below):

\begin{description}
    \item[\texttt{\textbackslash centering}] centres the content.
        It is the right command inside a float;
        the environment \texttt{center} adds extra space above and below, on top of the float's own.
    \item[\texttt{\textbackslash caption}] sets the caption, and with it the number.
        Without a caption there is no number, and nothing to refer to.
    \item[\texttt{\textbackslash label}] gives the float a name to refer to it (\autoref{sec:labels}).
        It must come \emph{after} \code|\caption|.
\end{description}

The order of caption and label matters because \code|\label| attaches to the last number that was set.
Before \code|\caption|, that is the number of the current section,
so a reference to the table shows the section number instead.
\hologo{LaTeX} gives no error and no warning, just a wrong number.

\section{Captions}\label{sec:captions}

Compiled with the rows filled in, the float from \autoref{sec:floats} looks like this:

\sliderender{04-table-float-render.pdf}

There is almost no gap between the caption and the table.
The standard classes assume captions go below, and only leave space above them.
The package \texttt{caption} fixes that, without any change in the document:

\slidecodetitled{04-caption-preamble.tex}{Preamble}
\sliderender{04-caption.pdf}

Its options change the look of all captions at once.
\texttt{labelfont=bf} sets \enquote{Table~1:} in bold,
\texttt{font=small} makes the whole caption smaller:

\slidecodetitled{04-caption-bf-preamble.tex}{Preamble}
\sliderender{04-caption-bf.pdf}

The convention is: tables have their caption above, figures below.
A table is read from the top, so its caption is read first;
a figure is looked at first, and the caption explains it.

\subsection{Lists of figures and tables}

\code|\listoffigures| and \code|\listoftables| work like \code|\tableofcontents|,
built from the captions.
A long caption would make a long entry;
the optional argument of \code|\caption| gives a short version for the list:

\owncode{ch06-lof.tex}
% trim = left bottom right top, in bp: the two lists and the figure
\ownrender[trim=110 490 110 118, clip]{ch06-lof.pdf}

Like the table of contents, the lists need two compilations (\autoref{sec:compiling}).

\section{Images}\label{sec:images}

Images are included with the package \texttt{graphicx}:

\slidecodetitled{04-graphicx.tex}{Preamble}

\begin{codeblock}
\includegraphics[width=0.8\textwidth]{figures/cooling-curves.pdf}
\end{codeblock}

The options go in the square brackets, the file in the braces.
\code|\includegraphics| works anywhere, even in the middle of a line;
for a figure with a caption and a number, it goes into a float (\autoref{sec:figures}).

\begin{itemize}
    \item \code|width=0.8\textwidth| scales the image to \qty{80}{\percent} of the text width.
        Relative to the text, it survives a change of margins or paper size; \code|width=12cm| does not.
        Other options: \code|height=5cm|, \code|scale=0.5|,
        \code|angle=90| (rotates),
        and \code|trim=1cm 0 1cm 0, clip| (cuts off the left and right \qty{1}{\centi\metre}).
        With \code|width| and \code|height| together, \code|keepaspectratio| keeps the image from being distorted.
    \item The path starts at the folder of the main \texttt{.tex} file.
        Keep images in a folder of their own, e.g.\ \texttt{figures/},
        use forward slashes, also on Windows,
        and avoid spaces in file names.
        \code|\graphicspath{{figures/}}| in the preamble would let you leave out the folder,
        but writing it out keeps it clear where a file is.
    \item Write the extension.
        \hologo{LaTeX} would find the file without it,
        but \texttt{figures/} might hold \texttt{cooling-curves.pdf} next to the \texttt{cooling-curves.tex} or \texttt{.py} that made it.
\end{itemize}

\subsection{File formats}

\hologo{pdfLaTeX} reads PDF, PNG and JPG:

\begin{description}
    \item[PDF] for vector graphics: plots and diagrams.
        They stay sharp at any zoom, and their text stays searchable.
    \item[PNG] for screenshots and other pixel images with sharp edges.
    \item[JPG] for photos.
\end{description}

SVG is not among them.
\hologo{pdfLaTeX} only reads formats that it can copy into the PDF more or less as they are,
and SVG would need a full renderer.
Export SVG images as PDF, e.g.\ with Inkscape (\textbf{Save As}, PDF),
or straight from the program that made them; plotting libraries can usually save PDF directly.
The package \texttt{svg} converts SVG files during compilation,
but only by running Inkscape, which needs \texttt{-shell-escape} and Inkscape installed.

\subsection{Figures}\label{sec:figures}

A figure is a float like a table, with \texttt{figure} instead of \texttt{table},
and \code|\includegraphics| instead of \texttt{tabular}.
The same \code|[htbp]|, \code|\centering|, \code|\caption| and \code|\label|, in the same order,
except that the caption comes below:

\slidecode{04-figure.tex}
\sliderender{04-figure.pdf}

\subsection{Several images side by side}

The package \texttt{subcaption} puts several images, each with its own caption, into one figure:

\owncodetitled{ch06-subfigures-preamble.tex}{Preamble}
\owncodetitled{ch06-subfigures.tex}{Document}
\ownrender{ch06-subfigures.pdf}

Each \texttt{subfigure} gets a width; the two widths plus the \code|\hfill| between them fill the line.
The sub-captions are labelled (a), (b),
and a reference to \texttt{fig:second} gives \enquote{1b}.
